\scanpage{129}
ce groupe $G=\pi/\pi'=\operatorname{Gal}(\mathbb K/\mathbb Q)$ opère continûment, et
il est engendré topologiquement par les $ {}^\sigma\dot\tau$, où $\dot\tau\in H$
est l'image de $\tau$, et $\sigma\in G$. Donc il \uline{est annulé par} $2$.

\margin{$\left[\rule{0pt}{9em}\right.$}Soit $L$ une extension abélienne de $\mathbb K$, \deleted{\unsure{envisag|correspondant} \unsure{les conditions}} \inserted{dont le}
groupe de Galois est donc un quotient $\Gamma$ de $H$; \deleted{Conditions}
Supposons $L\ne\mathbb K$, i.e. $\Gamma\ne\{e\}$. Conditions équivalentes

(i) les $\sigma\tau\sigma^{-1}$ $(\sigma\in\pi)$ coïncident toutes sur $L$
(donc y coïncident avec $\tau:\lambda\longrightarrow\bar\lambda$)

(ii) $L$ est une extension \uline{quadratique} de $\mathbb K$ \uline{galoisienne}
\uline{sur} $\mathbb Q$.

\deleted{\unsure{L est une extension} \illegible{2}} \inserted{(iii) $L$ est une extension \uline{quadratique} \uline{purement imaginaire} de $\mathbb K$.}

De plus, une telle $L$ \deleted{existe et} est unique, [et
d'ailleurs existe : il suffit de prendre une extension
quadratique imaginaire $E$ de $\mathbb Q$\inserted{(p. ex. $\mathbb Q(\sqrt{-1})$)}, et $L=\mathbb Q(\mathbb K,E)$
(p. ex. $L=\mathbb K(\sqrt{-1})$)].

\uline{Dém} $\mathrm{(i)}\Rightarrow\mathrm{(ii)}$ Comme $\Gamma$ est engendré par les images
\deleted{$\unsure{{}^\sigma\dot\tau}$} des $ {}^\sigma\tau$, \deleted{\unsure{il s'}} et que l'hyp. de (i) signifie que
ces images sont égales, on voit que (i) implique que
$\Gamma$ est engendré par l'élément $\dot\tau$ de degré $\mid 2$, donc est
$\simeq\mathbb Z/2\mathbb Z$ (car $\Gamma\ne(e)$ par hypothèse). D'ailleurs, si un
tel $\Gamma$ existe, \deleted{\unsure{on}} \inserted{il est} nécessairement égal à $H_G$ (car \deleted{\unsure{l'on}}
c'est un quotient de $H_G$, et l'argument précédent montre
que \deleted{\unsure{ce}} $H_G$ est engendré par un élément d'ordre $\mid2$. \deleted{\unsure{On}}
\deleted{il existe bel et} Il en résulte aussitôt que $L$ est galoisien
sur $\mathbb Q$, donc $\mathrm{(i)}\Rightarrow\mathrm{(ii)}$. D'autre part, $\mathrm{(ii)}\Rightarrow\mathrm{(i)}$, car
(ii) implique que $\Gamma\simeq\mathbb Z/2\mathbb Z$\inserted{ et que $G$ y opère}, et comme $\unsure{\Gamma|F}$ est engendré par les
images des $\sigma\tau\sigma^{-1}={}^{\sigma}\tau$, i.e. par les $ {}^\sigma\dot\tau$, on voit que
$\dot\tau$ (donc les $ {}^\sigma\dot\tau$) sont tous égaux à l'unique élément non nul de $\Gamma$.
$(\mathrm{(i)},\mathrm{(ii)})\Longleftrightarrow\mathrm{(iii)}$ car du moment que $L$ est quadratique sur $\mathbb K$,
la condition (i) sur les $\sigma\tau\sigma^{-1}$ est \unsure{évidemment|évidemt} équival\deleted{\unsure{ente}} à (iii)
% archive sheet number: 128
